\section{换帅背后：不同阶段的历史使命}

前面讲技术和组织，本章回到换帅。小鹏自动驾驶经历多任负责人：谷俊丽、吴新宙、李力耘、刘先明。刘先明没有直接定义每个人的历史使命，但可以从技术阶段看：早期需要探索和搭建能力，中期需要规则/量产体系，之后需要端到端转型，而现在进入 Physical AI 和更大模型/数据 scaling 阶段。

\begin{figure}[H]
\centering
\includegraphics[width=0.96\textwidth]{leader-mission-stages.png}
\caption{小鹏智驾负责人使命：不同阶段对应不同技术路线和组织任务。自制概念图，依据 00:54:30--01:48:46 对谈内容整理。}
\end{figure}

\begin{knowledgebox}{读图：换帅不是孤立人事，而是阶段变化}
每一任负责人面对的主问题不同：早期探索、规则量产、端到端转型、Physical AI 和下一阶段生活服务化。组织换帅往往伴随技术栈和公司战略变化。
\end{knowledgebox}

\subsection{何小鹏的角色}

刘先明描述何小鹏是坦诚、技术宅、问题很 sharp 的老板。他们最初聊的不是普通面试，而是“下一代怎样甩开对手”。在技术规划上，何小鹏会在机会与风险之间动态平衡；当模型性能和机会足够清晰时，他会调整判断。这个过程不是靠“说服老板”，而是靠把事情做出来，让技术涌现改变决策。

\begin{importantbox}{技术涌现改变决策}
在高不确定技术路线里，最有说服力的不是口头争论，而是模型性能突然跨过某个阈值。组织应该保留足够空间，让正确技术在指标和体验上自然显现。
\end{importantbox}

\subsection{预算、风险和 CEO 的问题}

何小鹏会问非常直接的问题，例如“为什么花这么多钱”。这个问题不是简单压预算，而是在逼团队说明：更大模型、更大数据、更大 infra 的投入，如何变成可部署能力、用户体验和商业结果。Physical AI 的 scaling 很贵，CEO 的角色不是只批准预算，而是持续追问机会、风险和时间窗口。

\begin{warningbox}{Scaling 不是免费午餐}
更大模型和更多数据会带来能力，但也会带来训练成本、车端部署成本、组织复杂度和质量风险。车企必须把技术路线和财务、量产、安全一起算。
\end{warningbox}

\subsection{Meta 与 Google 的组织启发}

刘先明提到 Meta 的半年节奏和 reorg 文化适合互联网快速试错，但 AI 时代很多事情需要从头建体系，半年一次的大调整未必友好。Google 在 Larry Page / Sergey Brin 回归后也发生变化。这个比较说明，公司文化会影响长线 AI 项目：如果考核周期太短，团队可能很难做需要一年以上才能验证的底层系统。

\begin{knowledgebox}{AI 时代的组织时间尺度}
互联网产品可以高频试错、快速砍掉；Physical AI 和大模型基础设施往往需要长周期验证。组织如果只接受短周期收益，就会伤害底层能力建设。
\end{knowledgebox}

\subsection{未来一到五年}

刘先明对一年、三年、五年的判断很有层次：一年后，智能化会更全面地接触生活，可能仍以数字 AI 为主；三年后，物理 AI 可能真正进入生活；五年后难以预测，因为世界发展太快。对个人和组织来说，能做的是持续学习，识别思维惯性，把当前认知下最正确的事情做好。

\subsection{本章小结}

换帅背后，是小鹏智驾从规则、端到端走向 Physical AI 的阶段切换。领导者的任务不只是管理团队，而是用工程和组织让正确技术路线跑出来。

